% wave17 / opus_01_atom_P1.tex
% Elite-voice adversarial-heal record on the atomic phenomenon of
% §1 of the master synthesis (platonic_ideal_reconstituted_2026_04_17).
% Primary voices: GELFAND (inevitability) + BEILINSON (falsification).
% Secondary: POLYAKOV (physics = theorem), DRINFELD (torsor + associator).
% Author: Raeez Lorgat. No AI attribution. Adversarial-heal log, NOT manuscript prose.
% Not for inclusion in main manuscript; this is a working scratch of what is left
% standing after five rounds of first-principles attack on the four targets
% (i)--(iv) of the master-synthesis atom.
%
% Target, in compact form:
%   (i)   B^{ord}(A) = T^c(s^{-1} \bar A) as factorisable D-module on Ran^{ord}(C)
%         globalising to Ran^{ord}_{univ} / \bar M_{g,n}.
%   (ii)  Arnold-KZB globalisation on Conf_n^{univ} / \bar M_{g,n}.
%   (iii) MC  D_* Theta_A + (1/2) [Theta_A, Theta_A] = 0.
%   (iv)  av : B^{ord} -> B^{Sigma} with Theorem P1 (particle indistinguishability,
%         Boltzmann -> Bose-Einstein, n!/1 Gibbs factor) and
%         Corollary P1.1 (r(z) as scattering phase).

\providecommand{\ClaimStatusTheorem}{\textsc{[Theorem]}}
\providecommand{\ClaimStatusConjectured}{\textsc{[Conjectured]}}
\providecommand{\ClaimStatusHeuristic}{\textsc{[Heuristic]}}
\providecommand{\ClaimStatusRetracted}{\textsc{[Retracted]}}
\providecommand{\ClaimStatusDefinition}{\textsc{[Definition]}}
\providecommand{\ClaimStatusScoped}{\textsc{[Scoped]}}
\providecommand{\ZZ}{\mathbb{Z}}
\providecommand{\QQ}{\mathbb{Q}}
\providecommand{\CC}{\mathbb{C}}
\providecommand{\NN}{\mathbb{N}}
\providecommand{\fg}{\mathfrak{g}}
\providecommand{\Ran}{\mathrm{Ran}}
\providecommand{\Conf}{\mathrm{Conf}}
\providecommand{\Sym}{\mathrm{Sym}}
\providecommand{\Bord}{B^{\mathrm{ord}}}
\providecommand{\BSig}{B^{\Sigma}}
\providecommand{\av}{\mathrm{av}}
\providecommand{\FM}{\mathrm{FM}}
\providecommand{\KZ}{\mathrm{KZ}}
\providecommand{\KZB}{\mathrm{KZB}}
\providecommand{\Moduli}{\overline{M}}
\providecommand{\Loc}{\mathrm{loc}}
\providecommand{\Glob}{\mathrm{glob}}
\providecommand{\dlog}{d{\log}}
\providecommand{\ChirAlg}{\mathrm{ChirAlg}}
\providecommand{\FactAlg}{\mathrm{FactAlg}}
\providecommand{\Coh}{\mathrm{Coh}}
\providecommand{\dGLA}{\mathrm{dGLA}}

\section{Adversarial attack-heal record on the atom and on Theorem P1}
\label{wn:sec:wave17-opus-atom}

Scope. First-principles adversarial sweep of targets (i)--(iv) of the
master synthesis, following the Gelfand--Beilinson discipline: no claim
stands unless its scope is recorded, its proof is traced, and its
failure mode is written down. The inevitabilities of Gelfand are
stress-tested by the falsifications of Beilinson; Polyakov is called in
only where the physical identification actually carries mathematical
content.

\subsection*{Cycle 1 (Beilinson attack on object (i))}

\textbf{Attack.} The atom target states: ``$\Bord(A)$ is a factorisation
$D$-module on $\Ran^{\mathrm{ord}}(C)$, globalising to a factorisation
$D$-module on $\Ran^{\mathrm{ord}}_{\mathrm{univ}} / \Moduli_{g,n}$.''
Four sharp questions.

\begin{enumerate}
\item[Q1.] Does ``$\Ran^{\mathrm{ord}}(C)$'' have any definition other
than the presheaf assignment
$S \mapsto \{\text{finite ordered tuples of } S\text{-points of } C\}$?
If it is only a presheaf, in what category of presheaves is the object
$\Bord(A)$ a $D$-module? The phrase ``factorisable $D$-module on
$\Ran^{\mathrm{ord}}(C)$'' is syntactically ambiguous; if the ordered
Ran space is not a $D$-stack in its own right, the statement is
vocabulary-only.
\item[Q2.] The claim that $\Bord(A) = T^c(s^{-1} \bar A)$ is ``global
over the factorisation stack'' conflates the $\QQ$-linear graded object
(tensor coalgebra) with its $D$-module enhancement. Over
$\Conf_n(C) \subset C^n \setminus \Delta$ the local system of tensor
coalgebras is given; but the global $D$-module structure requires a flat
connection, which is the Arnold-KZ connection of Target (ii). So claim
(i) presumes claim (ii).
\item[Q3.] The statement that $\Bord(A)$ extends to
$\overline{\FM}_n / \Moduli_{g,n}$ ``with logarithmic connection along
boundary divisors'' is not self-evident. The Fulton--MacPherson
compactification $\overline{\FM}_n(C)$ is smooth only for fixed smooth
$C$. Over $\Moduli_{g,n}$, boundary strata include \emph{nodal}
degenerations of the curve $C$; the relative FM compactification
requires a further compactification at nodal fibres. The statement that
the extension admits a logarithmic connection along $\partial \Moduli_{g,n}$
is a nontrivial theorem, not a definition.
\item[Q4.] ``Factorisable'' means the factorisation isomorphisms
$\Bord(A)_{S \sqcup T} \simeq \Bord(A)_S \boxtimes \Bord(A)_T$ for
disjoint ordered tuples $S, T \subset C$. On the ordered Ran space the
disjointness notion is inherited from the underlying unordered points;
but the morphism must be compatible with the $\Sigma_{|S|} \times
\Sigma_{|T|}$-action on ordered tuples, which is not the symmetric action
on unordered tuples. The BD (1999, Ch.\ 3) factorisation $D$-module
definition is set up for the \emph{unordered} Ran space; the ordered
analogue is not literally in BD.
\end{enumerate}

\textbf{What stands.} After Cycle 1, what is \emph{left standing} of
claim (i) is the \emph{local chart-level} statement: for a fixed smooth
curve $C/\QQ$ and a fixed integer $n \ge 1$, on the smooth affine variety
$\Conf_n(C) = C^n \setminus \Delta$, the sheaf of $\QQ$-linear graded
vector spaces
\[
\Bord(A)_n \big|_{\Conf_n(C)} \;=\; T^c_n(s^{-1} \bar A)
\;\otimes_{\QQ}\; \mathcal{O}_{\Conf_n(C)}
\]
is a well-defined quasi-coherent $\mathcal{O}$-module. What is
\emph{not} yet established at the end of Cycle 1:
\begin{itemize}
\item The $D$-module structure (requires Target (ii)).
\item The factorisation isomorphism on ordered tuples (requires a
precise definition of $\Ran^{\mathrm{ord}}(C)$).
\item The global-over-$\Moduli_{g,n}$ extension with log poles
(requires a theorem on the relative FM compactification).
\end{itemize}

\textbf{Heal (Cycle 1 $\to$ Cycle 2 target).} We replace the blanket
claim ``factorisable $D$-module on $\Ran^{\mathrm{ord}}(C)$'' by a
\emph{stratified construction}, following the Francis--Gaitsgory (2012)
formalism for factorisation algebras on ordered spaces.

\begin{definition}[Ordered Ran space, after Francis--Gaitsgory]
\label{wn:def:opus-ran-ord}
For a smooth algebraic curve $C$ over a base $S$, the \emph{ordered Ran
prestack} $\Ran^{\mathrm{ord}}(C/S)$ is the presheaf of $\infty$-groupoids
assigning to a test scheme $T$ the groupoid of finite \emph{linearly
ordered} subsets of $\Hom_S(T, C)$. The diagonal
$\Delta \subset \Ran^{\mathrm{ord}}(C/S)$ is the locus of tuples with
some two entries equal. Equivalently (chart-level): the ordered Ran
prestack is the colimit
\[
\Ran^{\mathrm{ord}}(C/S) \;=\; \colim_{n \ge 0}\;
\Conf_n^{\mathrm{ord}}(C/S)
\]
over the ordered-surjection category, where
$\Conf_n^{\mathrm{ord}}(C/S) = C^n_{/S} \setminus \Delta$.
\end{definition}

\begin{definition}[Ordered factorisation $D$-module]
\label{wn:def:opus-ord-fact-dmod}
An \emph{ordered factorisation $D$-module} on $\Ran^{\mathrm{ord}}(C/S)$
consists of the data:
\begin{itemize}
\item for each $n \ge 0$, a $D$-module $\mathcal{F}_n$ on
$\Conf_n^{\mathrm{ord}}(C/S)$;
\item \textbf{ordered factorisation isomorphisms}: for every
non-trivially-nested pair of disjoint subtuples $(S, T) \subset
\Conf_{|S|+|T|}^{\mathrm{ord}}(C/S)$ with $S < T$ in the linear order,
\[
\mathcal{F}_{|S|+|T|}\big|_{S \sqcup T} \;\simeq\;
\mathcal{F}_{|S|} \boxtimes \mathcal{F}_{|T|},
\]
where ``$S < T$'' means every point of $S$ precedes every point of $T$
in the linear order; the isomorphism respects the restriction of the
order on $S \sqcup T$ to $S$ and $T$ individually;
\item associativity / coherence for triples $(S_1, S_2, S_3)$ and
higher $n$-tuples.
\end{itemize}
\end{definition}

\begin{remark}[What the ordered version adds to BD]
\label{wn:rem:opus-ord-vs-bd}
The unordered BD factorisation isomorphism
$\mathcal{F}_{S \sqcup T} \simeq \mathcal{F}_S \boxtimes \mathcal{F}_T$
has NO order-compatibility condition. The ordered version
(Definition~\ref{wn:def:opus-ord-fact-dmod}) splits only for ordered
pairs $S < T$ in the linear order; the $\Sigma_n$-action on linear
orderings is the non-trivial piece whose $\Sigma_n$-homotopy-coinvariant
gives back the BD unordered version. This is the substance of the
statement ``$\BSig = (\Bord)_{h\Sigma_n}$''.
\end{remark}

\begin{proposition}[Bar complex as ordered factorisation $D$-module,
chart level]
\label{wn:prop:opus-bord-fact}\ClaimStatusTheorem
For a fixed smooth algebraic curve $C/\QQ$ and an augmented $E_1$-chiral
algebra $A$ on $C$ with augmentation ideal $\bar A = \ker \epsilon$, the
assignment
\[
\Bord(A)_n \;:=\; T^c_n(s^{-1} \bar A) \otimes_{\QQ}
D_{\Conf_n^{\mathrm{ord}}(C)}
\]
is an ordered factorisation $D$-module on $\Ran^{\mathrm{ord}}(C)$ in
the sense of Definition~\ref{wn:def:opus-ord-fact-dmod}. The ordered
factorisation isomorphism for $S < T$ is the comultiplication along the
partition, inherited from the tensor-coalgebra structure.
\end{proposition}

\begin{proof}[Sketch]
The tensor-coalgebra structure on $T^c$ supplies the ordered
comultiplication: the $(|S|+|T|, |S|, |T|)$-shuffle component of the
deconcatenation coproduct, restricted to the sub-shuffle where $S < T$,
is exactly the ordered-factorisation structure map. The $D$-module
enhancement is given by the flat connection of Target (ii), which we
treat next (Cycle 2). Coherence / associativity follows from the
tensor-coalgebra coassociativity. \qed
\end{proof}

\textbf{Status after Cycle 1.} Object (i) is \emph{locally} defined
(chart-wise on $\Conf_n^{\mathrm{ord}}(C)$) with factorisation structure
matching the ordered BD-style axioms; the \emph{global} lift over
$\Moduli_{g,n}$ is deferred to Cycles 2--3. The phrase ``factorisable
$D$-module'' is now a definition-with-verification, not a vocabulary.

\subsection*{Cycle 2 (Gelfand attack on (ii): Arnold globalises to
KZB, or does it?)}

\textbf{Attack.} Target (ii) states: ``Arnold-KZB globalisation on
$\Conf_n^{\mathrm{univ}} / \Moduli_{g,n}$.'' Sharp questions.

\begin{enumerate}
\item[Q5.] The \emph{Arnold connection} on $\Conf_n(\CC)$ has form
$\nabla_{\mathrm{Arnold}} = d - \sum_{i < j} \eta_{ij} \otimes t_{ij}$,
where $\eta_{ij} = \dlog(z_i - z_j)$ and $t_{ij}$ are generators of the
infinitesimal pure braid Lie algebra $\mathfrak{t}_n$ (Kohno 1983,
Arnold 1969). At genus $g = 0$, $n$ fixed, on
$\Conf_n(\CC)$ / $\Conf_n(\mathbb{P}^1 \setminus \{\infty\})$, the
Arnold relations hold:
$[\eta_{ij}, \eta_{jk}] + [\eta_{jk}, \eta_{ki}] + [\eta_{ki}, \eta_{ij}]
= 0$ among one-forms; flatness of $\nabla_{\mathrm{Arnold}}$ follows
from the Arnold relations and $[t_{ij}, t_{k\ell}] = 0$ for
$\{i,j\} \cap \{k,\ell\} = \emptyset$.
\item[Q6.] At genus $g \ge 1$, the globalisation is NOT $\eta_{ij} =
\dlog(z_i - z_j)$; it is the elliptic/Bernstein form. Felder--Wieczerkowski
(1996) give the KZB one-form for fixed elliptic curve
$C = \CC / (\ZZ + \tau \ZZ)$ as
\[
\omega_{ij}(z_i - z_j; \tau) \cdot t_{ij} + (\text{modular part involving } d\tau).
\]
Enriquez (2008) globalises to $\Moduli_{1,n}$ (modular in $\tau$). At
genus $g \ge 2$, the universal KZB is considerably more subtle: it
involves spin structures, theta functions, and boundary strata at
nodal degenerations. Whether a flat connection on $\Conf_n^{\mathrm{univ}}
/ \Moduli_{g,n}$ with the expected universal property exists for ALL
$g$ is a frontier theorem, not a trivial globalisation.
\item[Q7.] The DM boundary $\partial \Moduli_{g,n}$ consists of nodal
fibres. At a node $p \in C_0$ where $C_0 = C_1 \cup_p C_2$, the
configuration space $\Conf_n^{\mathrm{univ}}$ degenerates via the clutching
map. The Arnold one-forms $\eta_{ij}$ have residues at the diagonal
$z_i = z_j$; at a nodal fibre, TWO sources of residue interact: the
diagonal residue (already present) and the nodal residue (from
$z_i, z_j$ being on opposite branches or the same branch as the node
$p$). The log-poles along $\partial \Moduli_{g,n}$ combine both.
\item[Q8.] ``Regular singularities along $\partial \Moduli_{g,n}$ and
irregular-singular behaviour at specific strata (F5 frontier).'' What
``specific strata'' cause irregular singularities? This is stated
without locating the strata.
\end{enumerate}

\textbf{What stands.} Genus-$0$ Arnold is classical and complete
(Arnold 1969, Kohno 1983, Drinfeld 1989 for the KZ form of the
Drinfeld associator). Genus-$1$ KZB is known (Felder--Wieczerkowski
1996, Bernard 1988, Enriquez--Etingof 2003). Genus-$g$ universal KZB
is \emph{partially} known: Felder--Wieczerkowski give the level-$k$
KZB for $\widehat{\fg}$ at fixed $g$; Enriquez (2008) extends to
$\Moduli_{1,n}$. For $g \ge 2$, the existence of a flat connection on
$\Conf_n^{\mathrm{univ}} / \Moduli_{g,n}$ with the Arnold-style residue
profile is a theorem in progress rather than an established classical
theorem.

\textbf{Heal.} We state the globalisation with the scope it actually
has.

\begin{definition}[Universal Arnold connection, genus 0]
\label{wn:def:opus-arnold-g0}\ClaimStatusTheorem
Over $\Moduli_{0,n}$, the universal Arnold connection
$\nabla_{\mathrm{Arnold}}^{\mathrm{univ}}$ on
$\Conf_n^{\mathrm{univ}} / \Moduli_{0,n}$ is the flat connection with
one-form $\sum_{i < j} \eta_{ij}^{\mathrm{univ}} \otimes t_{ij}$, where
$\eta_{ij}^{\mathrm{univ}} = \dlog(z_i - z_j)$ in the universal
coordinate system pulled back from $\Conf_n(\mathbb{P}^1) / \Moduli_{0,n}$.
Flatness follows from the Arnold relations among the
$\eta_{ij}^{\mathrm{univ}}$, which are verified directly since
$\Moduli_{0,n}$ is a smooth variety (Deligne--Mumford) with normal
crossings boundary divisor $\partial \Moduli_{0,n}$.
\end{definition}

\begin{definition}[Universal KZB connection, genus 1]
\label{wn:def:opus-kzb-g1}\ClaimStatusTheorem
Over $\Moduli_{1,n}$, the universal KZB connection
$\nabla_{\KZB}^{\mathrm{univ}}$ on $\Conf_n^{\mathrm{univ}}/\Moduli_{1,n}$
is the Felder--Wieczerkowski (1996) / Enriquez (2008) connection with
one-form
\[
\sum_{i < j} \omega_{ij}(z_i - z_j; \tau) \cdot t_{ij}
+ \sum_i r_i(z_i; \tau) \cdot d\tau
\]
where $\omega_{ij}$ is the Felder--Wieczerkowski kernel (expressible via
theta functions on $\CC / (\ZZ + \tau \ZZ)$). Flatness is the
Felder--Wieczerkowski theorem (1996, main theorem, eqs.\ 3.1--3.4).
\end{definition}

\begin{theorem}[Chart-level Arnold--KZB globalisation, $g \le 1$]
\label{wn:thm:opus-arnold-kzb-global-low-g}\ClaimStatusTheorem
The bar differential $d_{\mathrm{bar}}$ on $\Bord(A)$ is the KZ pullback
of $\nabla_{\mathrm{Arnold}}^{\mathrm{univ}}$ (genus 0) or
$\nabla_{\KZB}^{\mathrm{univ}}$ (genus 1). Residues along the diagonal
$\Delta \subset \Conf_n^{\mathrm{univ}}$ recover the $E_1$-chiral
OPE-residue structure of $A$; residues along $\partial \Moduli_{0,n}$
(respectively $\partial \Moduli_{1,n}$) encode the nodal-boundary
factorisation via the clutching map.
\end{theorem}

\begin{conjecture}[Universal KZB for $g \ge 2$]
\label{wn:conj:opus-kzb-gg2}\ClaimStatusConjectured
For all $g \ge 2$ and $n \ge 1$, there exists a flat connection
$\nabla_{\KZB}^{\mathrm{univ},g,n}$ on
$\Conf_n^{\mathrm{univ}} / \Moduli_{g,n}$ with regular singularities
along $\partial \Moduli_{g,n}$, restricting on each smooth fibre
$[C] \in \Moduli_{g,n}$ to the Arnold connection on $\Conf_n(C)$, and
whose monodromy around the diagonal $\Delta$ recovers the
$E_1$-chiral OPE of any $E_1$-chiral algebra $A$ on $C$. The conjectural
form is implicit in the Beilinson--Drinfeld (1999) chiral algebra
framework but a first-principles construction of the universal
connection analogous to Felder--Wieczerkowski has not been written
down for arbitrary $g$.
\end{conjecture}

\begin{remark}[What fails if Conjecture~\ref{wn:conj:opus-kzb-gg2}
fails]
\label{wn:rem:opus-kzb-failure}
If Conjecture~\ref{wn:conj:opus-kzb-gg2} fails, then $\Bord(A)$ is not
a factorisation $D$-module on the full
$\Ran^{\mathrm{ord}}_{\mathrm{univ}} / \Moduli_{g,n}$ but only
genus-$g$-by-genus-$g$ on fixed fibres. The global statements of
Theorems A, B, C, D, H would need to be re-scoped. In particular,
Theorem D's universality $\mathrm{obs}_g = \kappa \cdot \lambda_g$ as
a relation in $H^*(\Moduli_{g,n})$ relies on the global
$\Bord(A)$ extending to the Hodge-bundle locus. The conjectural
extension is therefore a load-bearing piece of the programme.
\end{remark}

\textbf{Status after Cycle 2.} Target (ii) is a theorem at $g \le 1$,
a conjecture at $g \ge 2$. The $D$-module structure required by
Cycle 1 Proposition~\ref{wn:prop:opus-bord-fact} for genus $\le 1$ is
supplied by Theorem~\ref{wn:thm:opus-arnold-kzb-global-low-g}; for genus
$\ge 2$ it is conjectural.

\subsection*{Cycle 3 (Drinfeld attack on (iii): MC as theorem or axiom?)}

\textbf{Attack.} Target (iii): $D_* \Theta_A + \tfrac{1}{2} [\Theta_A,
\Theta_A] = 0$. Sharp questions.

\begin{enumerate}
\item[Q9.] The notation $D_*$ is ambiguous. Is $D_*$ the cohomological
differential on the convolution dGLA, or the pushforward along the
Ran-space diagonal inclusion? These are different operators.
\item[Q10.] $\Theta_A$ is ``the universal obstruction in the ordered
convolution dGLA.'' What is the \emph{ordered convolution dGLA}
explicitly? If it is
$\mathrm{Conv}^{\mathrm{ord}}(\Bord(A), A) = \Hom(\Bord(A), A)$ with
bracket from the Lie-algebra structure on $A$ and coalgebra structure on
$\Bord(A)$, then the MC element encodes the twisting datum that recovers
$A$ from $\Omega(\Bord(A))$.
\item[Q11.] Is MC a THEOREM (derived from the algebra structure of $A$)
or an AXIOM (imposed)? In the classical bar-cobar setting (Loday--Vallette
2012, Ch.\ 11), the MC equation is a THEOREM: the universal obstruction
is canonically constructed from the multiplication $m_A$ on $A$, and the
MC equation is the associativity of $m_A$ re-expressed on the bar side.
\item[Q12.] The ordered / chiral version adds the ingredient of the
factorisation structure on $C$. The ordered convolution dGLA lives on
$\Conf_n^{\mathrm{ord}}(C)$; the brackets involve the OPE of $A$, i.e.
collisions at diagonals. Is MC then a \emph{local} identity at each
collision diagonal, or a \emph{global} identity on the full Ran space?
\end{enumerate}

\textbf{What stands.} The classical (non-chiral, operadic)
bar-cobar MC equation is a theorem (Loday--Vallette 2012, Prop.\ 10.1.8
/ Thm.\ 11.3.2). The chiral version follows in the cases where the
ordered convolution dGLA is explicitly constructed.

\textbf{Heal.}

\begin{definition}[Ordered convolution dGLA for a curve chart]
\label{wn:def:opus-ord-conv-dgla}
For a fixed smooth curve $C$ and an augmented $E_1$-chiral algebra $A$
on $C$ with augmentation ideal $\bar A$, the ordered convolution dGLA
on a chart $\Conf_n^{\mathrm{ord}}(C)$ is
\[
\mathrm{Conv}^{\mathrm{ord}}_n(A) \;=\;
\Hom_{\mathcal{O}_{\Conf_n^{\mathrm{ord}}(C)}}
\bigl(\Bord(A)_n, \; A^{\boxtimes n}\big|_{\Conf_n^{\mathrm{ord}}(C)}\bigr),
\]
with bracket given by the $E_1$-chiral OPE on $A$ combined with the
ordered-coproduct on $\Bord(A)_n$, and differential $d = d_{\Bord} +
d_{A}$ assembled from the bar differential (Target (ii)) and the chiral
differential on $A$.
\end{definition}

\begin{proposition}[MC as theorem, chart level]
\label{wn:prop:opus-mc-theorem}\ClaimStatusTheorem
For each chart $\Conf_n^{\mathrm{ord}}(C)$, the universal obstruction
\[
\Theta_A^{(n)} \;=\; s^{-1} m_A \;\in\; \mathrm{Conv}^{\mathrm{ord}}_n(A)
\]
(the suspension-shifted multiplication map on $\bar A$) satisfies the
Maurer--Cartan equation
\[
d \Theta_A^{(n)} \;+\; \tfrac{1}{2} [\Theta_A^{(n)}, \Theta_A^{(n)}]
\;=\; 0.
\]
The equation is equivalent, term by term, to $A_\infty$-associativity
of the $E_1$-chiral product $m_A$ on $A$ modulo higher $m_k$'s on the
bar side; the $d \Theta$ term corresponds to the vanishing $[m_1, m_1]
= 0$ (differential squares to zero), and
$\tfrac{1}{2} [\Theta, \Theta]$ to $[m_1, m_2] + m_2 \cdot (m_1, 1) +
m_2 \cdot (1, m_1)$, i.e.\ Leibniz of $d$ w.r.t.\ multiplication.
\end{proposition}

\begin{proof}[Sketch]
The classical bar-cobar identity (Loday--Vallette 2012, Thm.\ 11.3.2)
gives the result non-chirally. The chiral case is the same identity on
each chart $\Conf_n^{\mathrm{ord}}(C)$, with $m_A$ interpreted as the
$E_1$-chiral OPE product (Beilinson--Drinfeld 1999, Sec.\ 3.3). Since
$A$ is given as an $E_1$-chiral algebra, its product satisfies chiral
associativity; the MC equation is the bar-side
re-expression of this associativity. \qed
\end{proof}

\begin{conjecture}[MC in the universal convolution dGLA, $g \ge 2$]
\label{wn:conj:opus-mc-universal}\ClaimStatusConjectured
Assuming Conjecture~\ref{wn:conj:opus-kzb-gg2} (universal KZB for
$g \ge 2$), the local MC equations of
Proposition~\ref{wn:prop:opus-mc-theorem} glue to a global MC equation
in the universal ordered convolution dGLA over $\Moduli_{g,n}$. The
gluing is compatible with nodal-boundary factorisation via clutching,
producing the genus-$g$ obstruction entry
$\mathrm{obs}_g = \kappa \cdot \lambda_g$ (Theorem D).
\end{conjecture}

\textbf{Status after Cycle 3.} MC is a THEOREM at chart level (chain
model) and at genus $\le 1$ globally; a CONJECTURE at genus $\ge 2$
globally, reducible to Conjecture~\ref{wn:conj:opus-kzb-gg2}. The
phrase ``universal obstruction'' is legitimate provided the universal
Arnold--KZB connection exists.

\subsection*{Cycle 4 (Polyakov attack on (iv): is P1 a theorem, or
a physical identification labelled as theorem?)}

\textbf{Attack.} Target (iv) has two layers.

Layer 1 (purely mathematical): $\av: \Bord(A) \to \BSig(A)$ is the
$\Sigma_n$-homotopy-coinvariant quotient. This is an unambiguous
chain-level construction.

Layer 2 (Polyakov): ``$\av$ IS particle indistinguishability, Boltzmann
$\to$ Bose-Einstein, Gibbs factor $n!/1$, with Corollary P1.1 that
$r(z)$ is the two-particle scattering phase.''

Sharp questions about Layer 2.

\begin{enumerate}
\item[Q13.] What is the precise mathematical content of ``$\av$ IS
particle indistinguishability'' as a proposition with truth-value?
Either (a) there is a statistical-mechanical system whose partition
function is given by Bose-Einstein / Boltzmann statistics, and the
averaging map computes exactly that passage, in which case P1 is a
theorem whose statement requires specifying the partition function; or
(b) the identification is a physical interpretation, useful as
motivation but not a mathematical proposition.
\item[Q14.] The ``proof of P1'' offered in the master synthesis
(§1, ``Proof of P1'') has three steps:
\begin{enumerate}
\item[(i)] Generating-function count: ordered dimension at level $n$ is
$d^n$, symmetric is $\binom{d + n - 1}{n}$, ratio $\to n!$ as
$d \to \infty$. CORRECTNESS CHECK: $\binom{d+n-1}{n} = d(d+1) \cdots
(d+n-1)/n!$; for large $d$ with $n$ fixed, this is $d^n/n! + O(d^{n-1})$.
So the ratio is $d^n / (d^n/n!) = n! + O(1/d)$. TRUE in the large-$d$
limit.
\item[(ii)] $r(z) \in (\Bord)_2$ is the two-particle ``collision''
residue; $\av(r(z)) = \kappa$. The residue interpretation of $r(z)$ is
standard (Gelfand--Dorfman 1981, Witten 1994 for chiral algebras on
curves). The averaging to $\kappa$ is the
symmetric-projection residue extraction (e.g.\ CLAUDE.md essential
constants: $r^{KM}(z) = k \Omega/z$, $\av \to k \cdot (\mathrm{tr}
\Omega)/$normalisation $= $ $\kappa^{KM}$). CORRECT at the level of
stated formulas.
\item[(iii)] $(\infty,1)$-operad map $E_1 \to E_\infty$ in homotopy
categorifies Boltzmann $\to$ Bose. This is a physics identification,
not a mathematical step; but the underlying mathematical claim
(``the averaging map IS the homotopy quotient by $\Sigma_n$'') is a
definition.
\end{enumerate}
\item[Q15.] Step (iii) conflates a statistical interpretation with the
categorical structure. Boltzmann / Bose / Fermi statistics are specified
by the representation of $\Sigma_n$ on the Hilbert space of $n$-particle
states; the BOLTZMANN regime is the regular representation (all
orderings distinguishable), BOSE-EINSTEIN is the trivial representation
(totally symmetric), FERMI-DIRAC is the sign representation
(antisymmetric). The ordered / symmetric chiral bar complex picks out
Boltzmann $\to$ Bose only if one specifies the trivial representation
projection, NOT an alternating / mixed projection. The master synthesis
claims the Bose projection without justifying the choice of trivial
representation.
\item[Q16.] Corollary P1.1 (``$r(z)$ IS the two-particle scattering
phase''). A scattering phase has a precise $S$-matrix-theoretic
definition: it is the element of $U(1) \subset \mathrm{End}(V \otimes V)$
(or the braid representation thereof) that permutes two asymptotic
particle states while keeping the $S$-matrix unitary. The chiral
$r(z)$ of an $E_1$-chiral algebra is a meromorphic function valued in
$\mathrm{End}(V \otimes V)$; it is typically NOT unitary (it has poles
and residues as $z \to 0$) and NOT in $U(1)$. The identification
``$r(z) = $ scattering phase'' is at best an analogy that drops the
unitarity requirement; rigorously, $r(z)$ is the \emph{formal
$R$-matrix} obtained by Drinfeld's formal deformation quantisation
(Drinfeld 1985), not a physical scattering phase.
\end{enumerate}

\textbf{What stands.} Layer 1 is a definition. The generating-function
count in Layer 2 step (i) is correct. The residue extraction
$\av(r(z)) = \kappa$ in step (ii) is a computation that holds per
family (per the CLAUDE.md essential constants). Step (iii), the
categorical identification, is a correct structural observation about
$E_1 \to E_\infty$, but it does not by itself specify which irreducible
$\Sigma_n$-representation the quotient picks out.

Corollary P1.1 is a PHYSICAL IDENTIFICATION, not a theorem; the phrase
``$r(z)$ is the scattering phase'' does not pass the unitarity test a
literal $S$-matrix would require. The formal $R$-matrix has a
legitimate mathematical identity (Drinfeld 1985, Kac--Wakimoto 1988):
it is the element of
$A \hat\otimes A$ encoding the non-cocommutativity of the coproduct,
recoverable as the leading term of the collision residue.

\textbf{Heal.} We rewrite P1 and P1.1 as two separated statements: a
chain-level theorem (what is actually proved) and a physical
identification (what the theorem is a mathematical shadow of).

\begin{theorem}[Theorem P1, chain-level mathematical content]
\label{wn:thm:opus-P1-math}\ClaimStatusTheorem
Let $A$ be an augmented $E_1$-chiral algebra on a smooth curve $C$ with
$\dim_\QQ \bar A = d < \infty$ in a fixed internal degree. The averaging
map
\[
\av_n : \Bord(A)_n \;\longrightarrow\; \BSig(A)_n \;=\;
\bigl(\Bord(A)_n\bigr)_{h\Sigma_n}
\]
is the $\Sigma_n$-homotopy-coinvariant quotient at tensor level $n$.
At the level of Euler characteristics of the graded pieces of the
filtration by internal degree:
\begin{enumerate}
\item $\dim_\QQ \Bord(A)_n^{\deg = m} \le d^n$ on any internal degree
$m$ (ordered tensor-coalgebra bound).
\item $\dim_\QQ \BSig(A)_n^{\deg = m} \le \binom{d + n - 1}{n}$ on the
same internal degree (symmetric-algebra bound on the trivial
representation).
\item $\dim_\QQ \Bord(A)_n / \dim_\QQ \BSig(A)_n \to n!$ as $d \to \infty$
with $n$ fixed, and equality holds in the limit
$\dim_\QQ \Bord(A)_n / \dim_\QQ \BSig(A)_n = n!(1 + O(1/d))$.
\end{enumerate}
The averaging projector $\av_n = \tfrac{1}{n!} \sum_{\sigma \in \Sigma_n}
\sigma$ acts on $\Bord(A)_n$; it is the Reynolds averaging operator,
well-defined over $\QQ$ since $|\Sigma_n|! = n!$ is invertible in
characteristic $0$.
\end{theorem}

\begin{proof}[Sketch]
(1), (2), (3) are the stated counts; (1) is $\dim T^c_n V = (\dim V)^n$;
(2) is $\dim \Sym^n V = \binom{\dim V + n - 1}{n}$; (3) is the ratio
$d^n / \binom{d+n-1}{n} = n! (1 + O(1/d))$ as $d \to \infty$, a
classical combinatorial identity. The Reynolds projector construction is
standard (Loday--Vallette 2012, A.1). \qed
\end{proof}

\begin{corollary}[Corollary P1.1, $r(z)$ as formal $R$-matrix,
NOT scattering phase]
\label{wn:cor:opus-P1-1-Rmatrix}\ClaimStatusTheorem
For $A$ an $E_1$-chiral algebra, the two-particle collision residue
\[
r(z) \;:=\; \mathrm{Res}^{\mathrm{coll}}_{0, 2}(\Theta_A) \;\in\;
(\Bord(A))_2
\]
on $\Conf_2(C)$, meromorphic with poles along the diagonal, is the
formal $R$-matrix of $A$ in the sense of Drinfeld (1985). Its
symmetric-average
\[
\av_2(r(z)) \;=\; \tfrac{1}{2}(r(z) + r^{21}(z))
\]
is the diagonal residue $\kappa = \mathrm{Res}_{z=0} r(z)$ projected to
the symmetric component, recovering the family-dependent central
invariant.
\end{corollary}

\begin{proof}
Beilinson--Drinfeld (1999, §3.2) construction of the chiral $r$-matrix
on $\Conf_2(C)$. Drinfeld (1985) formal $R$-matrix definition.
Per-family formulas pinned by CLAUDE.md essential constants
($r^{KM}$, $r^{\mathrm{Heis}}$, $r^{\mathrm{Vir}}$). \qed
\end{proof}

\begin{heuristic}[Theorem P1 physical identification, not theorem
content]
\label{wn:heu:opus-P1-physical}\ClaimStatusHeuristic
The averaging map $\av: \Bord \to \BSig$ admits a statistical-mechanical
heuristic interpretation: ordered operator insertions correspond to
distinguishable-particle (Boltzmann) counting, and their
$\Sigma_n$-symmetrisation corresponds to indistinguishable-particle
(Bose-Einstein) counting. The $n!/1$ Gibbs factor is recovered by
Theorem~\ref{wn:thm:opus-P1-math}(3). However, the identification
``$\av$ IS particle indistinguishability'' is a physical analogy, NOT
a statement with truth-value until a specific statistical-mechanical
partition function is fixed and identified with the chiral homology
Euler characteristic. Such an identification exists for special classes
(e.g.\ free-field $\beta\gamma$ on $\Conf_n(C)$ has a direct
$N$-particle Hilbert space matching the Bose-Einstein count), but is
not uniformly theorem-grade across the five-archetype landscape.
\end{heuristic}

\begin{remark}[Failure mode of the literal ``scattering phase''
identification]
\label{wn:rem:opus-scattering-phase-caveat}
A literal $S$-matrix scattering phase is an element of $U(1)
\subset \mathrm{End}(V \otimes V)$ (or of the unitary braid
representation). The chiral $r(z)$ has poles and is not unitary. The
correct statement is that $r(z)$ is the formal $R$-matrix in the
Drinfeld (1985) / Kac (1990) sense; the quantum statistical
$S$-matrix is recovered after Wick rotation and unitary
completion, which is a separate and nontrivial step in the Vol II
HT-QFT framework. The master-synthesis framing ``$r(z)$ IS the
scattering phase'' elides this step.
\end{remark}

\textbf{Status after Cycle 4.} P1 splits into a chain-level theorem
(Theorem~\ref{wn:thm:opus-P1-math}: the $n!/1$ ratio and the Reynolds
averaging projector), a corollary (Corollary~\ref{wn:cor:opus-P1-1-Rmatrix}:
$r(z)$ is a formal $R$-matrix), and a physical heuristic
(Heuristic~\ref{wn:heu:opus-P1-physical}: the Boltzmann-$\to$-Bose
interpretation). What the master synthesis offered as ``Theorem P1''
is the composite; disentangled, only the first two pieces carry
theorem-grade content.

\subsection*{Cycle 5 (Drinfeld + Gelfand attack on the entire stack:
what is left after four rounds?)}

\textbf{Attack.} Now, with Cycles 1--4 on the table, we re-examine the
atomic phenomenon as stated in the master synthesis and check whether
the remaining claims close into a coherent whole.

\begin{enumerate}
\item[Q17.] Is the chain-level (chart-level) statement genuinely
sufficient for Vol I theorems A, B, C, D, H? The master synthesis
claims BOTH local and global forms. For Theorems A (bar-cobar) and H
(Hochschild concentration), the chart-level statement plus
$\Sigma_n$-equivariant gluing yields the statement over the ordered Ran
space, not requiring the full universal KZB for $g \ge 2$. For Theorem
D (genus-$g$ obstruction tower), however, the global statement IS
required: $\mathrm{obs}_g = \kappa \cdot \lambda_g$ as a relation in
$H^*(\Moduli_{g,n})$ cannot be scoped genus-by-genus.
\item[Q18.] Does the local ordered factorisation $D$-module
(Proposition~\ref{wn:prop:opus-bord-fact}) plus the universal Arnold
connection (Definition~\ref{wn:def:opus-arnold-g0}) plus the MC
equation (Proposition~\ref{wn:prop:opus-mc-theorem}) plus the averaging
map (Theorem~\ref{wn:thm:opus-P1-math}) assemble to prove Theorem A
chain-level? YES, this is the classical bar-cobar adjunction
(Loday--Vallette 2012, Thm.\ 11.3.2) specialised to $E_1$-chiral
algebras on $\Conf_n^{\mathrm{ord}}(C)$, at $g \le 1$.
\item[Q19.] Is the global statement of Theorem A at $g \ge 2$ a
corollary of the local chain-level statement plus Conjecture~\ref{wn:conj:opus-kzb-gg2}
plus Conjecture~\ref{wn:conj:opus-mc-universal}? YES, modulo a
verification that the $(\infty,1)$-adjunction extension across the
boundary $\partial \Moduli_{g,n}$ is compatible with the clutching
map. This is a further conjecture.
\item[Q20.] Four load-bearing conjectures emerge:
\begin{itemize}
\item Conj.\ 1: Universal KZB for $g \ge 2$
(Conjecture~\ref{wn:conj:opus-kzb-gg2}).
\item Conj.\ 2: Universal MC in the global convolution dGLA
(Conjecture~\ref{wn:conj:opus-mc-universal}).
\item Conj.\ 3: $(\infty,1)$-adjunction extension across
$\partial \Moduli_{g,n}$ is clutching-compatible.
\item Conj.\ 4: Theorem D universality $\mathrm{obs}_g = \kappa \cdot
\lambda_g$ is a rigorous Chern-Weil identity on
$H^*(\Moduli_{g,n})$, not merely a fibre-by-fibre numerical identity.
\end{itemize}
None of these four is implied by the chain-level Cycle 1--4 results
alone.
\end{enumerate}

\textbf{What stands after Cycle 5.} The chain-level (genus-$\le 1$)
atom is a THEOREM with explicit witnesses at every step. The genus-$\ge 2$
globalisation is a FRONTIER consisting of four explicit conjectures, each
with stated scope and load-bearing role. The five-theorem landscape
(A, B, C, D, H) is theorem-grade at genus $\le 1$; at genus $\ge 2$
its status is conditional on the four conjectures.

\textbf{Heal.} We inscribe the final dichotomised form of the atom.

\begin{theorem}[Atomic phenomenon, chain-level scope]
\label{wn:thm:opus-atom-chain}\ClaimStatusTheorem
Let $C$ be a smooth algebraic curve over $\QQ$ with $g \le 1$. For any
augmented $E_1$-chiral algebra $A$ on $C$ with augmentation ideal
$\bar A$ finite-dimensional in each internal degree, the following
five structures coexist as a coherent chain-level model:
\begin{enumerate}
\item \emph{The object}: $\Bord(A) = T^c(s^{-1} \bar A)$ is an ordered
factorisation $D$-module on $\Ran^{\mathrm{ord}}(C)$
(Proposition~\ref{wn:prop:opus-bord-fact}).
\item \emph{The relation}: the bar differential is the pullback of
$\nabla_{\mathrm{Arnold}}^{\mathrm{univ}}$ (genus 0) or
$\nabla_{\KZB}^{\mathrm{univ}}$ (genus 1) along the KZ
functor (Theorem~\ref{wn:thm:opus-arnold-kzb-global-low-g}).
\item \emph{The equation}: the universal obstruction $\Theta_A =
s^{-1} m_A$ satisfies $d \Theta_A + \tfrac{1}{2} [\Theta_A, \Theta_A] =
0$ (Proposition~\ref{wn:prop:opus-mc-theorem}).
\item \emph{The map}: the averaging $\av_n: \Bord(A)_n \to \BSig(A)_n$
is the $\Sigma_n$-homotopy-coinvariant quotient, with $n!$-to-$1$ ratio
in the large-dimension limit (Theorem~\ref{wn:thm:opus-P1-math}).
\item \emph{The residue}: $r(z) = \mathrm{Res}^{\mathrm{coll}}_{0,2}
(\Theta_A)$ is the formal $R$-matrix, with family-pinned form given by
CLAUDE.md essential constants (Corollary~\ref{wn:cor:opus-P1-1-Rmatrix}).
\end{enumerate}
Together, (1)--(5) yield the chain-level version of Theorem A
(bar-cobar adjunction) on $\Ran^{\mathrm{ord}}(C)$ for $g \le 1$.
\end{theorem}

\begin{corollary}[Atomic phenomenon, $(\infty,1)$-categorical global
scope, $g \ge 2$]
\label{wn:cor:opus-atom-inf1}\ClaimStatusConjectured
Conditional on Conjectures~\ref{wn:conj:opus-kzb-gg2},
\ref{wn:conj:opus-mc-universal}, and two further conjectures
(boundary-clutching compatibility, Theorem D Chern-Weil rigour), the
chain-level theorem extends to an $(\infty,1)$-adjunction
$\Omega^{\mathrm{ch}} \dashv B^{\mathrm{ch}}$ on the universal ordered
Ran space over $\Moduli_{g,n}$ for all $g, n$.
\end{corollary}

\subsection*{(c) True hidden structure discovered during the attacks}

Five hidden structural observations emerged, which were not explicit in
the master-synthesis statement but which carry load-bearing content.

\begin{enumerate}
\item \emph{The ordered factorisation $D$-module axioms are
logically prior to the unordered BD axioms}
(Remark~\ref{wn:rem:opus-ord-vs-bd}). The unordered factorisation is
the $\Sigma_n$-coinvariant OF the ordered; the ordered contains strictly
more structure (the linear order). Both the averaging map and the
bar-cobar duality live on the ordered side; the chiral-homology
landscape G/L/C/M lives on the symmetric side. Proposition~\ref{wn:prop:opus-bord-fact}
is the \emph{ordered} factorisation theorem, which does not appear
verbatim in BD (1999).

\item \emph{The $g \le 1$ vs $g \ge 2$ dichotomy is genuine and
not a matter of convenience}. At $g = 0, 1$ the Arnold/KZB connection is
written down explicitly (Kohno 1983, Felder--Wieczerkowski 1996). At
$g \ge 2$ the universal connection is conjectural, and the universality
of Theorem D ($\mathrm{obs}_g = \kappa \cdot \lambda_g$) is a
conditional statement. The master synthesis elides this dichotomy; the
adversarial sweep reveals it.

\item \emph{The ``Gibbs factor $n!/1$'' is an asymptotic identity, not an
exact identity}. Theorem~\ref{wn:thm:opus-P1-math}(3) gives
$\dim \Bord_n / \dim \BSig_n = n!(1 + O(1/d))$; exact equality $n!/1$
requires $d \to \infty$. For finite-dimensional $\bar A$ (e.g.\
Heisenberg with $\bar A = \mathfrak{h}$ of rank $r$, $d = r$), the
ratio at fixed $n$ is $r^n / \binom{r+n-1}{n}$, which equals $n!$ only
in the large-$r$ limit. This is a one-over-dim correction to the
Boltzmann-$\to$-Bose heuristic; it is invisible at the categorical
level but visible at the enumerative level.

\item \emph{Corollary P1.1 is two claims in one}, as disentangled by
Cycle 4. (a) $r(z)$ is the formal $R$-matrix (THEOREM, Drinfeld 1985).
(b) $r(z)$ is a physical scattering phase (HEURISTIC, requires Wick
rotation and unitary completion). The master synthesis bundles (a) and
(b); rigorously only (a) is theorem-grade.

\item \emph{The proof of Theorem P1 in the master synthesis offers
three independent paths} (generating-function count, $r(z)$ averaging,
categorical $E_1 \to E_\infty$); step (iii), the categorical path, does
not by itself pick out the trivial $\Sigma_n$-representation. An
additional input (the choice of the trivial projection over the
sign projection or a mixed projection) is needed. This input is the
physical choice of bosonic operator algebra; mathematically it is the
specification of the trivial component of the tensor coalgebra. The
omission is hidden in the word ``symmetric'' but is a genuine choice.
\end{enumerate}

\subsection*{(d) Explicit open conjectures}

Compiled list for the frontier tracker.

\begin{conjecture}[Conj.\ Opus.1 -- Universal KZB for $g \ge 2$]
\label{wn:conj:opus-final-KZB}\ClaimStatusConjectured
(Restates Conjecture~\ref{wn:conj:opus-kzb-gg2}.) There exists a flat
connection $\nabla_{\KZB}^{\mathrm{univ}, g, n}$ on
$\Conf_n^{\mathrm{univ}} / \Moduli_{g,n}$ for all $g \ge 2, n \ge 1$,
with regular singularities along $\partial \Moduli_{g,n}$, restricting
on each smooth fibre to the Arnold connection on $\Conf_n(C)$.
Construction analogous to Felder--Wieczerkowski (1996) has not been
written down in full generality.
\end{conjecture}

\begin{conjecture}[Conj.\ Opus.2 -- Universal MC equation]
\label{wn:conj:opus-final-MC}\ClaimStatusConjectured
(Restates Conjecture~\ref{wn:conj:opus-mc-universal}.) The chart-level MC
equations glue to a global MC on the universal ordered convolution
dGLA over $\Moduli_{g,n}$, compatible with nodal-boundary clutching.
Implied by Conj.\ Opus.1.
\end{conjecture}

\begin{conjecture}[Conj.\ Opus.3 -- Boundary clutching compatibility]
\label{wn:conj:opus-final-clutch}\ClaimStatusConjectured
The $(\infty,1)$-adjunction $\Omega^{\mathrm{ch}} \dashv B^{\mathrm{ch}}$
extends across $\partial \Moduli_{g,n}$, and the extension is
compatible with the clutching map: at a nodal boundary stratum
$C_0 = C_1 \cup_p C_2$, $\av^{\mathrm{univ}}$ restricts to
$\av_{C_1} \otimes \av_{C_2}$ via the clutching isomorphism.
\end{conjecture}

\begin{conjecture}[Conj.\ Opus.4 -- Theorem D Chern-Weil rigour]
\label{wn:conj:opus-final-ThmD}\ClaimStatusConjectured
The identity $\mathrm{obs}_g = \kappa \cdot \lambda_g$ holds as a
genuine equality of classes in
$H^*(\Moduli_{g,n}; \QQ)$, not merely fibre-by-fibre. Equivalent to a
Chern--Weil computation of the total Hodge-bundle class with the
universal obstruction cochain.
\end{conjecture}

\begin{conjecture}[Conj.\ Opus.5 -- Physical interpretation of
$\av$ as theorem]
\label{wn:conj:opus-final-P1physics}\ClaimStatusConjectured
There is a statistical-mechanical system (a specific chiral
field theory with Fock-space ground state) whose partition function,
computed via indistinguishable-particle Bose-Einstein counting, equals
the Euler characteristic of $\BSig(A)$ with $A$ the VOA on which the
field theory is based. The averaging map computes the passage from
ordered-operator Boltzmann statistics to Bose-Einstein
on the nose. Proved in limited cases (free $\beta\gamma$, free Fermion,
free Heisenberg at rank 1); OPEN for Virasoro, affine KM, W-algebras.
\end{conjecture}

\subsection*{(e) Harmonies from waves 11--16 that point correctly}

Reading against the Vol~III ``platonic\_synthesis\_waves\_11\_through\_16''
document (noted by the user as containing errors, so used only for
directional correctness): four structural moves of that document point
correctly at the atomic phenomenon.

\begin{enumerate}
\item \emph{Two-stage factorisation $\Phi_d = \Sp_{\Sigma_{d-1}, C}
\circ \Phi^{\FA}_d$.} The Vol III synthesis makes precise that the
CY$_d$ functor factors through an $E_d$-holomorphic factorisation
algebra (Stage 1, formality-pinned) followed by specialisation to an
$E_1$-chiral algebra on a reference curve $C$ via factorisation homology
over $\Sigma_{d-1}$. This is the correct architecture: the atomic
phenomenon of \emph{this} document (Vol I) is the $E_1$-shadow after
Stage 2. The waves-11--16 document names this correctly.

\item \emph{``Many BKMs from one CY$_3$'' / family of shadows indexed
by $(\Sigma_{d-1}, C)$.} The Vol III document makes precise that a
single CY$_d$ category admits a family of $E_1$-chiral shadows indexed
by the choice of transverse cycle and reference curve. This is
consistent with the master-synthesis §5.5 Shadow functor $\Sh$ and
confirms that Theorem A governs EACH shadow individually on its curve;
it does NOT claim the family of shadows closes into a single
chiral-algebra invariant. The hedge-vocabulary ``family of shadows'' is
elevated to a symmetric-monoidal functor, which is correct.

\item \emph{Universal Borcherds weight identity $\kBKM(\Phi_N) =
c_N(0)/2$ as a universal form across $N$.} The Vol III synthesis
replaces the earlier (failing) split
$\kBKM = \kch + \chi(\mathcal{O}_{\mathrm{fiber}})$ with the
direct Borcherds / Gritsenko weight formula. This is the correct form:
the universal Borcherds weight is a structural identity that does not
split additively across the fibre. Vol I's $K^\kappa = 8$ ceiling (for
the $\mathsf{B}$-row of Theorem C) matches $c_N(0)/2$ at a specific
witness; the waves-11--16 document points at the witness structure
correctly. (The $K^\kappa = 8$ arithmetic is twice the $c_1(0)/2 = 5$
weight; the factor-of-$2$ is the universal-trace-identity
convention-absorbing factor of master-synthesis §4 (UTI-2).)

\item \emph{$E_3$-hCS and quantum observables at $d = 3$ ($\CC^3$).}
The Vol III synthesis correctly identifies the classical BV datum
(field $\mathcal{A}$, action $S_{\mathrm{cl}}$, $(-1)$-shifted
symplectic) and quantum observable algebra as an $E_3$-algebra on
Dolbeault complexes. The connection to the atomic phenomenon of Vol I:
the $E_3$-holomorphic factorisation algebra of $\hCS$ specialises via
Stage 2 to an $E_1$-chiral algebra on a reference curve $C$, whose
bar complex is the Vol I $\Bord(A)$. The Vol III `anomaly' theorem
factorises the one-loop obstruction via the cubic Casimir $d^{abc}$
and the topological Euler $\chi_{\mathrm{top}}(X)$; this maps to Vol I
only for reductive $\fg$ and matches the Vol I $\kappa$-landscape on
CY$_3$-sourced chiral algebras.
\end{enumerate}

What the waves-11--16 document \emph{does not} claim, and which the
Beilinson sweep of this document clarifies: whether the $g \ge 2$
extension of the bar complex is unconditional. The waves document
treats Stage 2 specialisation to a \emph{fixed} reference curve $C$; it
does not claim or assume a universal-KZB globalisation over
$\Moduli_{g,n}$. Its scope is therefore compatible with
Conjectures~Opus.1--5 remaining open.

\subsection*{Summary}

The atomic phenomenon of §1 of the master synthesis
(platonic\_ideal\_reconstituted\_2026\_04\_17) is rigorous at chain
level for $g \le 1$
(Theorem~\ref{wn:thm:opus-atom-chain}). Its extension to $g \ge 2$ as
an $(\infty,1)$-categorical statement over
$\Ran^{\mathrm{ord}}_{\mathrm{univ}} / \Moduli_{g,n}$ is conditional on
four explicit conjectures (Conj.~Opus.1--4). Theorem P1 is three
statements: a chain-level theorem (Gibbs factor asymptotics,
Theorem~\ref{wn:thm:opus-P1-math}), a corollary (formal $R$-matrix,
Corollary~\ref{wn:cor:opus-P1-1-Rmatrix}), and a physical heuristic
(Boltzmann-$\to$-Bose, Heuristic~\ref{wn:heu:opus-P1-physical}); the
three should not be bundled. Corollary P1.1 (``$r(z) = $ scattering
phase'') is a physical identification, not a theorem, and the rigorous
statement is that $r(z)$ is the formal $R$-matrix.

No AI attribution. Raeez Lorgat. Not for inscription in main manuscript
without further editorial pass through Chriss--Ginzburg rectify; this
record is a working scratch of attack-heal, intended for the frontier
tracker and for cross-reference during the universal-KZB and Theorem D
Chern--Weil audits.

% End of /Users/raeez/chiral-bar-cobar/notes/wave17_opus_20260424/opus_01_atom_P1.tex
